\chapter{Architettura e implementazione}\label{ch:sistema}
\section{Separazione fra osservazione, azione e valutazione}
L'architettura separa stato osservabile, policy e ambiente offline per rendere verificabile il confine fra informazione storica e valutazione futura. Distinguere generazione dei candidati, contesto e ranking permette di cambiare policy mantenendo lo stesso spazio d'azione e le stesse osservazioni ammissibili; isolare reward e coda differita impedisce che un feedback calcolabile dall'evaluatore diventi subito apprendimento. La modularità serve quindi causalità e confrontabilità, oltre alla manutenzione del codice.

Il loader conserva tutte le righe in una cache SQLite temporanea ordinata cronologicamente. Stato e candidati si aggiornano incrementalmente; indici e spool limitano la RAM occupata dai record. Sono risposte tecniche al volume dei dati, non alternative architetturali confrontate scientificamente.

Il flusso logico è:
\begin{center}
\fbox{\parbox{.91\linewidth}{\centering
TSV e identità canoniche $\longrightarrow$ eventi cronologici\\[4pt]
$D_{<t}\longrightarrow$ TemporalState $\longrightarrow$ CandidateGenerator\\[4pt]
$\longrightarrow$ ContextBuilder $\longrightarrow$ scaler congelato $\longrightarrow$ policy\\[4pt]
$\longrightarrow$ Top-1 e Top-K $\longrightarrow$ log immutabile e metriche\\[4pt]
Solo Top-1 $\longrightarrow$ reward offline $\longrightarrow$ coda a $t+7d$
$\longrightarrow$ aggiornamento della policy}}
\end{center}
L'accesso al futuro appartiene al calcolatore di reward e al valutatore, non al percorso che costruisce stato, candidati o feature. Ogni policy riceve lo stesso ambiente e lo stesso criterio di eleggibilità. Per ciascuna configurazione vengono costruite istanze indipendenti; non sono condivisi stato mutabile, feedback o flussi casuali.

\section{Stato incrementale e causalità}
Lo stato al decision point è
\begin{equation}
O_t=f(D_{<t}),\qquad D_{<t}=\{e:\operatorname{timestamp}(e)<t\}.\label{eq:state}
\end{equation}
TemporalState conserva conteggi per utente, per coppia utente--item e utente--artista, ultimo timestamp individuale, ascolti globali e utenti distinti per item e artista, oltre alla prima osservazione dell'item. Queste strutture supportano novelty, recency e age senza filtrare nuovamente l'intero dataset per ciascuna decisione.

Gli indici del CandidateGenerator ricevono gli eventi tramite listener solo dopo il loro ingresso nello stato. Per ciascun gruppo valutabile di decisioni simultanee, \metric{run_phase()} applica nell'ordine: feedback con scadenza $\leq t$, avanzamento dello stato a $t$, tutte le decisioni del gruppo e infine ammissione degli eventi a $t$. Candidati e contesti condividono dunque $D_{<t}$; nessuna scelta del gruppo genera un feedback immediato.

La memoria di apprendimento ha una convenzione distinta dallo stato delle feature. Se una decisione precedente è presa a $s$ e $s+7d=t$, la sua reward include gli ascolti esattamente a $t$ e può aggiornare $A,b$ prima delle nuove scelte. Per esempio, un ascolto dell'artista selezionato alle 10:00 dell'8 gennaio conta nel feedback della decisione delle 10:00 del 1 gennaio, ma non nei conteggi o nei candidati delle decisioni delle 10:00 dell'8 gennaio. Il simulatore dà priorità ai feedback maturi rispetto alle decisioni simultanee: non usa ascolti successivi a $t$, ma l'intera informazione del learner non è descritta dal solo $D_{<t}$. È la convenzione implementata per l'uguaglianza dei timestamp, non una disponibilità anticipata rispetto alla scadenza.

Per esempio, se la prima sessione contiene eventi alle 10:00, 10:03 e 10:05 e una nuova sessione inizia alle 10:20, la decisione delle 10:20 osserva i tre eventi precedenti. L'evento delle 10:20 rimane invisibile, anche a un secondo utente che decida nello stesso istante. La costruzione di contesti usa esattamente lo stato avanzato a quel timestamp e rifiuta item con prima osservazione non anteriore a $t$.

\section{Tre pool e unione dei candidati}
Personal dà accesso agli interessi già osservati, Popular aggiunge segnali globali e Discovery ammette item mai ascoltati dall'utente. La loro unione evita di vincolare tutte le policy alla sola storia individuale e definisce un ambiente comune: il generatore stabilisce quali azioni siano disponibili, la policy come ordinarle. La scoperta rimane possibile solo entro questo catalogo osservato.

Per ogni utente si costruisce
\begin{equation}
\mathcal A_u(t)=A_N^{\mathrm{personal}}(u,t)\cup
A_N^{\mathrm{popular}}(t)\cup A_N^{\mathrm{discovery}}(u,t).\label{eq:pool}
\end{equation}
$N\in\{50,100,200\}$ è la profondità di \emph{ciascuna} sorgente, non la cardinalità finale. Dopo la deduplicazione, $|\mathcal A_u(t)|\leq3N$; l'overlap può renderla sensibilmente inferiore. Non si impongono quote fra sorgenti nella lista finale.

Il pool Personal prende i primi $N$ item, in ordine crescente di chiave canonica, fra le tracce osservate degli artisti già conosciuti: comprende così item già ascoltati e item non ascoltati di quegli artisti, senza priorità obbligatoria al primo gruppo. Popular prende i primi $N$ per ascolti globali dell'item prima di $t$, con tie-break sulla chiave. Discovery prende i primi $N$ item mai ascoltati dall'utente nello stesso ordine lessicografico e li etichetta \emph{familiar artist} ($C_{u,i}=0,C_{u,g}>0$) o \emph{new artist} ($C_{u,i}=C_{u,g}=0$). L'etichetta non modifica l'estrazione: non bilancia le due categorie né impone loro quote. La regola è deterministica, ma introduce un possibile bias lessicografico; il generatore non campiona uniformemente il catalogo.

Le sorgenti vengono troncate separatamente, quindi unite nell'ordine Personal, Popular, Discovery. Un duplicato conserva tutti i source flag. L'ordinamento dei pool non coincide con quello della policy: in particolare, Popular usa popolarità \emph{track-level}, mentre la baseline Global Popularity usa popolarità \emph{artist-level}. Tale differenza è intenzionale nel codice verificato.

\section{Vettore contestuale e standardizzazione}\label{sec:features}
Il contesto associa attività dell'utente, ampiezza e ripetizione dei gusti a popolarità globale, familiarità utente--item/artista, tempo e novelty. Questa rappresentazione permette a LinUCB di combinare segnali che le baseline usano separatamente o ignorano, generalizzando fra item. È una scelta compatta e interpretabile; non descrive similarità audio o semantica. La valutazione riguarda l'insieme delle feature, senza attribuire vantaggi a singole componenti non sottoposte ad ablation.

Indichiamo con $E_u$ gli eventi storici, $D_{u,I}$ e $D_{u,G}$ gli item e gli artisti distinti, $P_i,P_g$ gli ascolti globali e $U_i,U_g$ gli utenti globali distinti. Tutte le quantità sono misurate prima di $t$. Poniamo $\ell(z)=\ln(1+z)$. La Tabella~\ref{tab:features} riporta l'intero schema, evitando l'equivoco fra dimensione totale e numero di feature standardizzate.

\begin{table}[H]\centering\small
\begin{tabular}{clp{8.4cm}}
\toprule Pos. & Espressione & Significato\\\midrule
1 & $\ell(E_u)$ & Attività cumulata dell'utente\\
2 & $\ell(D_{u,G})$ & Ampiezza storica degli artisti\\
3 & $\ell(D_{u,I})$ & Ampiezza storica degli item\\
4 & $1-D_{u,G}/E_u$ & Ripetizione artista, zero se $E_u=0$\\
5 & $1-D_{u,I}/E_u$ & Ripetizione item, zero se $E_u=0$\\
6 & $\ell(P_i)$ & Popolarità globale dell'item\\
7 & $\ell(U_i)$ & Numero di utenti distinti dell'item\\
8 & $\ell(P_g)$ & Popolarità globale dell'artista\\
9 & $\ell(U_g)$ & Numero di utenti distinti dell'artista\\
10 & $\ell(t-t_i^{\mathrm{first}})$ & Età dalla prima osservazione globale, in secondi\\
11 & $\ell(C_{u,i})$ & Conteggio storico utente--item\\
12 & $\ell(t-t_{u,i}^{\mathrm{last}})$ & Recency in secondi; zero se item mai ascoltato\\
13 & $\ell(C_{u,g})$ & Conteggio storico utente--artista\\
14 & $\mathbb 1[C_{u,i}=0]$ & Track novelty binaria\\
15 & $\mathbb 1[C_{u,g}=0]$ & Artist novelty binaria\\
\bottomrule
\end{tabular}
\caption{Le 15 componenti del ContextBuilder. Le prime 13 sono standardizzate; le ultime due restano binarie.}\label{tab:features}
\end{table}

Per $j\leq13$, lo scaler usa media e deviazione standard di popolazione calcolate sui contesti di warm-up:
\begin{equation}
z_j=(x_j-\mu_j)/\sigma_j,\qquad
\sigma_j^2=M^{-1}\sum_{m=1}^{M}(x_{m,j}-\mu_j)^2.
\end{equation}
La stima streaming segue Welford; una deviazione nulla viene sostituita da uno. I flag finali hanno media convenzionale zero e scala uno. Non sono introdotti un bias esplicito, feature audio, generi o embedding. In particolare non esiste una feature aggiuntiva di recency utente--artista, anche se lo stato conserva l'ultimo timestamp dell'artista.

Il campione di fitting conta ogni occorrenza candidato--decisione, per tutti gli utenti del warm-up, inclusi quelli non eleggibili in validation/test. Gli item presenti più spesso nei pool pesano quindi più volte. I contesti sono costruiti cronologicamente e poi aggregati dallo scaler: nessuno è ricostruito usando lo stato finale del warm-up.

\section{Definizione del reward e feedback differito}
In assenza di risposte a raccomandazioni registrate, gli ascolti futuri forniscono un proxy di compatibilità con il comportamento osservato. L'obiettivo artist-level riconosce anche l'ascolto di un brano diverso dello stesso artista; perde però precisione sul titolo, senza essere universalmente preferibile a un obiettivo track-level. La sola Top-1 identifica un'azione cui associare feedback, separandola dalla valutazione della lista. Sette giorni costituiscono un compromesso operativo fra risposta prossima e opportunità di riascolto, orientato dall'EDA preliminare ma non dimostrato ottimale. Il logaritmo attenua i conteggi elevati, mantenendo la distinzione fra nessun ascolto e ascolti ripetuti.

Il reward dell'azione $(u,i,t)$ dipende dall'artista $g(i)$:
\begin{equation}
r(u,i,t)=\ln\!\left(1+\#\{e:u(e)=u,\ g(e)=g(i),\
t<\tau(e)\leq t+\Delta\}\right),\quad \Delta=7d.\label{eq:reward}
\end{equation}
L'evento a $t$ è escluso, quello a $t+\Delta$ incluso. Gli indici futuri usano timestamp compatti e ricerche binarie. La reward può essere calcolata dall'ambiente dopo la scelta, ma non è fornita al learner prima della scadenza. Le baseline ignorano l'aggiornamento di reward; la loro storia continua a evolvere con gli ascolti reali.

\clearpage
\section{Warm-up, ottimizzazioni e costo computazionale}
Il fitting è indipendente dalla policy a parità di $N$. Il runner riutilizza dunque lo scaler per le configurazioni della stessa profondità, riproducendo in nuove istanze lo stato e gli indici fino al confine storico. Nel test $N=100$ il warm-up registrato conta 1.281.669 decision point e 372.870.964 contesti: la dimensione spiega perché non sia un'operazione banale. Il riuso riduce i fit da 36 a 3 in validation e a uno fra le cinque policy del test; non condivide la memoria di apprendimento.

La Tabella~\ref{tab:runtime} riporta durate effettive dei run held-out, senza nuovi benchmark. Per Random, prima configurazione, la colonna warm-up/replay include il costo del fitting comune e della preparazione; per le successive include il ripristino da parametri già verificati. Il totale per run comprende anche operazioni ulteriori rispetto alla sola simulazione, per esempio valutazione e serializzazione. Non rappresenta il tempo complessivo del processo con caricamento iniziale ed eleggibilità.

\begin{table}[htbp]\centering\small
\begin{tabular}{lrrr}
\toprule
Policy & Totale [min] & Warm-up/replay [min] & Simulazione [min] \\
\midrule
Random & 274,84 & 248,05 & 26,59 \\
Global Popularity & 30,78 & 4,79 & 25,81 \\
Personalized Greedy & 31,30 & 4,73 & 26,38 \\
$\varepsilon$-Greedy & 33,31 & 6,04 & 27,07 \\
LinUCB & 51,91 & 4,93 & 46,77 \\
\bottomrule
\end{tabular}

\caption{Durate ufficiali del test, in minuti. La prima riga sostiene il costo comune del warm-up e non va interpretata come lentezza della policy Random.}\label{tab:runtime}
\end{table}

Il costo dominante della prima configurazione deriva dalla preparazione di centinaia di milioni di contesti, non dal suo ranking casuale. Le successive istanze ricostruiscono la stessa storia in pochi minuti, senza rigenerare sessioni, pool e contesti del warm-up. La simulazione di LinUCB richiede più tempo delle baseline perché valuta lo score contestuale per ogni candidato e aggiorna la matrice del modello. Queste misure sono specifiche dell'esecuzione registrata: non isolano l'effetto di hardware, carico del sistema e I/O.

La cache SQLite è temporanea e la connessione viene chiusa prima della rimozione. I risultati conclusi sono invece persistenti. Tale separazione permette di conservare le evidenze anche dopo il cleanup e di riprendere una run interrotta verificando commit, configurazione, scaler e checksum dei dettagli. L'ottimizzazione riguarda la ripetizione del lavoro storico; il protocollo statistico e la sequenza dei feedback rimangono quelli del percorso di riferimento.
